\documentclass{math_library}

\begin{document}
\section{Теорема Гана-Банаха}

\begin{theorem}
Нехай $X$ --- нормований простір та $X_0 \subset X$ --- підпростір. Тоді для кожного функціонала $l \in X_0'$ існує продовження $\tilde{l} \in X'$ функціонала $l$, причому $\| \tilde{l} \| = \|l\|$.
\end{theorem}

Доведення доволі громіздке, тому розбив на підрозділи.

\subsubsection{Випадок, коли $X$ --- сепарабельний простір над $\mathbb{R}$}
Доведення цього пункту розіб'ємо ще на пару частин, почнемо з першої частини, що треба показати.
\begin{enumerate}[wide=0pt,label={\roman*.}]
\item \textit{$l$ можна продовжити на деякий підпростір $W$, де $X \supset W \supset X_0$.}\\
Нехай $X_0$ --- підпростір $X$ та $X_0 \neq X$. Зафіксуємо точку $y \notin X_0$ та розглянемо підпростір 
\[
W = \linspan\{X_0 \cup \{y\}\},
\]
тобто кожен елемент $z \in W$ записується як
\[
z = x + \lambda y, \qquad x \in X_0, \lambda \in \mathbb{R}.
\]
Визначимо функціонал $\tilde{l} \colon W \to \mathbb{R}$ ось таким чином: 
\[
\tilde{l}(x) = l(x) + \lambda c, \qquad c = \tilde{l}(y).
\]
За побудовою такий оператор --- лінійний. Залишилося підібрати таке $c \in \mathbb{R}$, щоб виконувалась рівність $\|\tilde{l}\| = \|l\|$ --- тим самим ми й обмеженість доведемо автоматично. Але згідно зі зауваження, нам треба підібрати $c \in \mathbb{R}$, щоб $\|\tilde{l}\| \leq \|l\|$.
\begin{itemize}[wide=0pt]
\item \textit{Випадок при $\lambda > 0$}.\\
Зафіксуємо $h_1,h_2 \in X_0$ та зауважимо, що справедлива нерівність:
\begin{align*}
l(h_2) - l(h_1) & = l(h_2-h_1) \leq |l(h_2-h_1)| \leq \|l\| \|h_2 - h_1\| \\
& = \|h \| \|(h_2+y) - (y+h_1)\| \leq \|l\| \|h_1+y\| + \|l\| \|h_2+y\|.
\end{align*}
Звідси випливає, що
\[
-\|l\| \|h_1+y\| - l(h_1) \leq \|l\| \|h_2+y\| - l(h_2).
\]
Оскільки це виконується для всіх $h_1,h_2 \in X_0$, то тоді
\[
a_1 = \sup_{h_1 \in X_0} (-\|l\| \|h_1+y\| - l(h_1)) \leq \inf_{h_2 \in X_0} ( \|l\| \|h_2+y\| - l(h_2)) = a_2.
\]
Оберемо число $c \in \mathbb{R}$ так, щоб $a_1 \leq c \leq a_2$. Звідси справедлива така нерівність для всіх $h \in X_0$:
\[
-\|l\| \|h+y\| - l(h) \leq c \leq \|l\| \|h+y\| - l(h).
\]
Тепер покладемо елемент $h = \lambda^{-1}g$ та домножимо обидві частини нерівності на $\lambda$. Оскільки ми домовилися $\lambda > 0$, то знаки нерівностей зберігаються. Коротше, отримаємо:
\begin{align*}
-\|l\| \|g+\lambda y \| - l(g) \leq \lambda c \leq \|l\| \|g + \lambda y \| - l(g); \\
-\|l\| \|g+\lambda y\| \leq l(g) + \lambda c \leq \|l\| \|g+\lambda y\|; \\
|\tilde{l}(x)| = |l(g) + \lambda c| \leq \|l\| \|g + \lambda y\| = \|l\| \|x\|.
\end{align*}
Власне, далі аналогічними міркуваннями (\textit{як в попередньому твердженні}) отримаємо $\|\tilde{l}\| \leq \|l\|$.

\item \textit{Випадок при $\lambda < 0$}.\\
Перепишемо елемент ось так:
\[
z = -(-x + (-\lambda)y).
\]
У нас тепер $-\lambda > 0$ та $-z = t = -x + (-\lambda)y$, звідси отримаємо
\[
|\tilde{l}(t)| \leq \|l\| \|t\| \implies |\tilde{l}(x)| \leq \|l\| \|x\|.
\]
\end{itemize}

\item \textit{Продовежння на нашому конкретному $X$ існує}.\\
Оскільки $X$ --- сепарабельний, то існує множина (\textit{ми оберемо зліченну}) $A = \{x_1,x_2,\dots\}$, яка є щільною підмножиною $X$. Також ми досі маємо підпростір $X_0 \subset X$.\\
Позначимо $x_{n_1} \in A$ --- перший з елементів, де $x_{n_1} \notin X_0$. За кроком i. існує продовження $l_1$ оператора $l$ на підпростір $W_1 = \linspan\{X_0 \cup \{x_{n_1}\}\}$.\\
Позначимо $x_{n_2} \in A$ --- перший з елементів, де $x_{n_2} \notin W_1$. За кроком i. існує продовження $l_2$ оператора $l_1$ на $W_2 = \linspan\{W_1 \cup \{x_{n_2}\}\}$.\\
\vdots \\
Отримаємо ланцюг підпросторів $X_0 \subset W_1 \subset W_2 \subset \dots$ та послідовність функціоналів $(l_n)_{n \in \mathbb{N}}$, для яких:
\begin{center}
$\forall n \geq 1: l_n \in W_n',\qquad {l_n}|_{X_0} = l, \qquad \|l_n\| = \|l\|$.
\end{center}
Покладемо множину, яка є лінійною:
\[
M = \displaystyle\bigcup_{n=1}^\infty W_n.
\]
Визначимо функціонал $L_0 \colon M \to \mathbb{R}$ ось так: 
\[
x \in M \implies x \in W_N \implies L_0(x) = l_N(x).
\]
Зрозуміло цілком, що $L_0$ --- лінійний, а також $\|L_0\| = \|l\|$. Оскільки $M \supset A$ та $A$ всюди щільна, то $M$ --- всюди щільна. Отже, за попереднім твердженням існує продовження $L \colon X \to \mathbb{R}$, для якого $\|L\| = \|L_0\| = \|l\|$.
\bigskip \\
Висновок: ми довели теорему Гана-Банаха для випадку, коли $X$ --- дійсний, сепарабельний простір.
\end{enumerate}

\subsubsection{Випадок, коли $X$ --- довільний нормований простір над $\mathbb{R}$}
Все ще $X_0 \subset X$. Позначимо за $l_p$ --- продовження $l$  зі збереженням норми на множині $P \supset X_0$. Таке продовження існує (\textit{див I. та i.}). Позначимо $E$ за множину всіх таких продовжень. На ній введемо відношення $\preceq$ таким чином:
\[
l_p \preceq l_q \iff P \subset Q \wedge \forall x \in P: l_Q(x) = l_P(x).
\]
Зрозуміло, що $\preceq$ задає відношення порядку, внаслідок чого $E$ --- частково впорядкована. Зафіксуємо $Y = \{l_{P_\alpha} \mid \alpha \in A\}$ --- будь-яку лінійно впорядкувану підмножину $X$. Знайдемо верхню грань.\\
Для цього покладемо $P_* = \displaystyle\bigcup_{\alpha \in A} P_\alpha$ та на множині $P_*$ задамо функціонал $l_*$ таким чином:\[
x \in P_* \implies x \in P_{\alpha_0} \implies l_*(x) = l_{\alpha_0}(x).
\]
Зрозуміло, що $l_*$ --- лінійний, причому $\|l_*\| = \|l\|$. На множині $\bar{P_*}$ продовжимо функціонал, як було в твердженні --- отримаємо функціонал $l_{\bar{P_*}}$, причому $\|l_{\bar{P_*}}\| = \|l_*\| = \|l\|$. Даний функціонал $l_{\bar{P}_*}$ на $\bar{P}_*$ буде верхньою гранню $Y$. Отже, за лемою Цорна, існує максимальний елемент $E$. Це буде функціонал $L$, який визначений на $X$ (\textit{у протилежному випадку його можна було би ще продовжити та він не був би максимальним елементом}).
\bigskip \\
Висновок: ми довели теорему Гана-Банаха для випадку, коли $X$ --- будь-який дійсний простір.

\subsubsection{Випадок, коли $X$ --- довільний нормований простір над $\mathbb{C}$}
Нехай $X$ --- комплексний лінійний нормований простір. Розглянемо одночасно $X_{\mathbb{R}}$ --- асоційований з $X$ дійсний нормований простір; тобто під час множення на скаляр ми допускаємо лише дійсні коефіцієнти. Зауважимо, що $X_{\mathbb{R}} = X$ як множини, утім не як простори.\\
Розглянемо довільний функціонал $l \colon X \to \mathbb{C}$. Раз $l(x) \in \mathbb{C}$, то для кожного $x \in X$ можна записати функціонал як $l(x) = m(x) + i n(x)$. У цьому випадку $m(x) = \Re l(x),\ n(x) = \Im l(x)$.

\begin{proposition}
Нехай $l \in X'$. Тоді $m,n \in (X_{\mathbb{R}})'$.
\end{proposition}

\begin{proof}
Нехай $\alpha,\beta \in \mathbb{R}$ та $x,y \in X$. Тоді ми отримаємо з одного боку
\[
l(\alpha x + \beta y) = m(\alpha x + \beta y) + i n(\alpha x + \beta y),
\]
а з іншого боку,
\[
l(\alpha x + \beta y) = \alpha l(x) + \beta l(y) = \alpha (m(x) + i n(x)) + \beta (m(y) + i n(y)) = (\alpha m(x) + \beta m(y)) + i(\alpha n(x) + \beta n(y)).
\]
Знаючи, що комплексне число рівне тоді й лише тоді, коли дійсні та уявні частини збігаються, отримаємо
\begin{align*}
m(\alpha x + \beta y) = \alpha m(x) + \beta m(y), \qquad n(\alpha x + \beta y) = \alpha n(x) + \beta n(y).
\end{align*}
Отже, $m,n$ --- лінійний функціонали. Обмеженість $m$ (\textit{аналогічно з $n$}) випливає з такго ланцюга нерівностей:
\[
|m(x)| \leq |m(x) + in(x)| = |l(x)| \leq \|l\| \|x\|. \qedhere
\]
\end{proof}

\begin{proposition}
$n(x) = -m(ix)$.\\
(\textit{іншими словами, ми можемо функціонал $l$ відновити повністю, знаючи функціонал $m$})
\end{proposition}

\begin{proof}
Для цього зауважимо наступне:
\[
m(ix) + in(ix) = l(ix) = i l(x) = i(m(x) + in(x)) = -n(x) + im(x) \implies n(x) = -m(ix).
\]
Отже, одержимо $l(x) = m(x) -im(ix)$.
\end{proof}

Тепер повертаємось до доведення теореми Гана-Банаха. Маємо $X \supset X_0$ --- два комплексних простори та $X_{\mathbb{R}}, {X_0}_{\mathbb{R}}$ --- асоційовані простори. Маємо функціонал $l \colon X_0 \to \mathbb{C}$, який визначається дійсним функціоналом $m \colon {X_0}_{\mathbb{R}} \to \mathbb{R}$. Оскільки це дійсний функціонал, ми можемо продовжити до $M \colon X_{\mathbb{R}} \to \mathbb{R}$ зі збереженням норми. Покладемо 
\[
L(x) = M(x) - i M(ix).
\]
Неважко буде довести, що $L$ --- комплексний лінійний функціонал. Залишилося довести, що $\|L\| = \|l\|$. Знову ж таки, достатньо довести $\|L\| \leq \|l\|$. Запишемо
\[
L(x) = |L(x)|e^{i \varphi},
\]
де $\varphi = \arg L(x)$. Тоді
\[
|L(x)| = e^{-i \varphi} L(x) = L(e^{-i \varphi} x) = M(e^{-i \varphi} x) = |M(e^{-i \varphi}x)| \leq \|M\| \|e^{-i \varphi} x\| = \|m\| \|x\| \leq \|l\| \|x\|.
\]
Отже, $\|L\|$. Ми тут юзали той факт, що $L(y) = M(y)$ при $L(y) \in \mathbb{R}$.

\begin{remark}
Зауважимо, що якщо $X_0$ --- лінійна множина (\textit{але не підпростір}), то теорема Гана-Банаха все одно виконується. У цьому випадку $\bar{X_0}$ буде підпростором $X$. Функціонал $l$ продовжується неперервним чином на $\bar{X_0}$, а далі застосовується доведена теорема.
\end{remark}

\iffalse
%спочатку запишу теорему для окремого випадку
\begin{theorem}[Теорема Гана-Банаха]
Нехай $E$ --- нормований простір та $G \subset E$ --- підпростір. Тоді для кожного функціонала $l \in G'$ існує продовження $\tilde{l} \in E'$ функціонала $l$, причому $\| \tilde{l} \| = \|l\|$.
\end{theorem}
\noindent
(TODO: форматувати текст доведення та переглянути всі моменти)
\fi

\iffalse
\begin{proof}
1. Обмежимось випадком, коли $E$ --- дісний та сепарабельний простір.\\
I. \textit{Доведемо, що $l$ можна продовжити на деякий підпростір $E \supset F \supsetneq G$.}\\
Нехай $G$ --- підпростір $E$ та $G \neq E$. Зафіксуємо $y \notin G$ та розглянемо підпростір $F = \linspan\{G \cup \{y\}\}$. Тобто кожний елемент $x \in F$ записується як $x = g + \lambda y$ при $g \in G, \lambda \in \mathbb{R}$. Визначимо оператор $\tilde{l}(x) = l(g) + \lambda c$, де $c = \tilde{l}(y)$. За побудовою, такий оператор --- лінійний.\\
Тепер залишилося підібрати таке $c \in \mathbb{R}$, щоб виконувалося $\|\tilde{l}\| = \|l\|$ --- тим самим ми й обмеженість доведемо автоматично. Але згідно зі зауваження, нам треба підібрати $c \in \mathbb{R}$, щоб $\|\tilde{l}\| \leq \|l\|$.\\
Обмежимося поки що $\lambda > 0$. Нехай зафіксовано $h_1,h_2 \in G$ та зауважимо, що справедлива нерівність:\\
$l(h_2) - l(h_1) = l(h_2-h_1) \leq |l(h_2-h_1)| \leq \|l\| \|h_2 - h_1\| = \|h \| \|(h_2+y) - (y+h_1)\| \leq \|l\| \|h_1+y\| + \|l\| \|h_2+y\|$.\\
Звідси випливає, що $-\|l\| \|h_1+y\| - l(h_1) \leq \|l\| \|h_2+y\| - l(h_2)$.\\
Оскільки це $\forall h_1,h_2 \in G$, то тоді $\displaystyle\sup_{h_1 \in G} (-\|l\| \|h_1+y\| - l(h_1)) \leq \inf_{h_2 \in G} ( \|l\| \|h_2+y\| - l(h_2))$.\\
Для зручності супремум позначу за $a_1$ та інфімум за $a_2$. Оберемо число $c \in \mathbb{R}$ так, щоб $a_1 \leq c \leq a_2$. Звідси справедлива така нерівність:\\
$\forall h \in G: -\|l\| \|h+y\| - l(h) \leq c \leq \|l\| \|h+y\| - l(h)$.\\
Тепер покладемо елемент $h = \lambda^{-1}g$ та домножимо обидві частини нерівності на $\lambda$. Оскільки ми домовилися $\lambda > 0$, то знаки нерівностей зберігаються. Коротше, отримаємо:\\
$-\|l\| \|g+\lambda y \| - l(g) \leq \lambda c \leq \|l\| \|g + \lambda y \| - l(g)$.\\
$-\|l\| \|g+\lambda y\| \leq l(g) + \lambda c \leq \|l\| \|g+\lambda y\|$.\\
$|\tilde{l}(x)| = |l(g) + \lambda c| \leq \|l\| \|g + \lambda y\| = \|l\| \|x\|$.\\
Власне, далі аналогічними міркуваннями (як в попередньому твердженні) отримаємо $\|\tilde{l}\| \leq \|l\|$.\\
Тепер що робити при $\lambda < 0$. Перепишемо $x = -(-g + (-\lambda)y)$. У нас тепер $-\lambda > 0$ та $-x = t = -g + (-\lambda)y$, звідси отримаємо\\
$|\tilde{l}(t)| \leq \|l\| \|t\| \implies |\tilde{l}(x)| \leq \|l\| \|x\|$.
\bigskip \\
II. \textit{Тепер доведемо, що продовежння на нашому конкретному $E$ існує}.\\
Оскільки $E$ --- сепарабельний, то існує (ми оберемо зліченну) множина $A = \{x_1,x_2,\dots\}$, яка є щільною підмножиною $E$. Також ми досі маємо $G \subset E$ --- підпростір.\\
Позначимо $x_{n_1} \in A$ --- перший з елементів, де $x_{n_1} \notin G$. За кроком І, існує $l_1$ --- продовження $l$ на $G_1 = \linspan\{G \cup \{x_{n_1}\}\}$.\\
Позначимо $x_{n_2} \in A$ --- перший з елементів, де $x_{n_2} \notin G_1$. За кроком І, існує $l_2$ --- продовження $l_1$ на $G_2 = \linspan\{G_1 \cup \{x_{n_2}\}\}$.\\
\vdots \\
Отримаємо ланцюг підпросторів $G \subset G_1 \subset G_2 \subset \dots$ та набір функціоналів $l_1,l_2,\dots$, для яких:\\
$\forall n \geq 1:$ \qquad $l_n \colon G_n \to \mathbb{R}$ --- обмежена; \qquad $l_n|_G = l$; \qquad $\|l_n\| = \|l\|$.\\
Покладемо множину $M = \displaystyle\bigcup_{n=1}^\infty G_n$, яка є лінійною. Визначимо функціонал $L_0 \colon M \to \mathbb{R}$ таким чином: $x \in M \implies x \in G_N \implies L_0(x) = l_N(x)$. Зрозуміло цілком, що $L_0$ --- лінійний, а також $\|L_0\| = \|l\|$.\\
Оскільки $M \supset A$ та $A$ всюди щільна, то $M$ --- всюди щільна. Отже, за попереднім твердженням, існує продовження $L \colon E \to \mathbb{R}$, для якого $\|L\| = \|L_0\| = \|l\|$.\\
Висновок: ми довели теорему Гана-Банаха для випадку, коли $E$ --- дійсний сепарабельний.
\bigskip \\
2. Тепер будемо доводити теорему для $E$ --- довільний дійсний нормований простір. Все ще $G \subset E$. Позначимо за $l_p$ --- продовження $l$  зі збереженням норми на множині $P \supset G$. Таке продовження існує див (1. та I.). Позначимо $X$ --- множина всіх таких продовжень. На ній введемо відношення $\preceq$ таким чином:\\
$l_p \preceq l_q \iff P \subset Q$ та $l_Q(x) = l_P(x), \forall x \in P$.\\
Зрозуміло, що $\preceq$ задає відношення порядку, внаслідок чого $X$ --- частково впорядкована. Зафіксуємо $Y = \{l_{P_\alpha} \mid \alpha \in A\}$ --- будь-яку лінійно впорядкувану підмножину $X$. Знайдемо верхню грань.\\
Для цього покладемо $P_* = \displaystyle\bigcup_{\alpha \in A} P_\alpha$ та на множині $P_*$ задамо функціонал $l_*$ таким чином:\\
$x \in P_* \implies x \in P_{\alpha_0} \implies l_*(x) = l_{\alpha_0}(x)$.\\
Зрозуміло, що $l_*$ --- лінійний, причому $\|l_*\| = \|l\|$. На множині $\bar{P_*}$ продовжимо функціонал, як було в твердженні --- отримаємо функціонал $l_{\bar{P_*}}$, причому $\|l_{\bar{P_*}}\| = \|l_*\| = \|l\|$. Даний функціонал $l_{\bar{P}_*}$ на $\bar{P}_*$ буде верхньою гранню $Y$. Отже, за лемою Цорна, існує максимальний елемент $X$. Це буде функціонал $L$, який визначений на $E$ (у протилежному випадку його можна було би ще продовжити та він не був би максимальним елементом).\\
Висновок: ми довели теорему Гана-Банаха для випадку, коли $E$ --- дійсний (не обов'язково сепарабельний) нормований простір.
\end{proof}
\noindent
Насправді, на цьому теорема Гана-Банаха ще не закінчена. Ми можемо її довести на випадок, коли нормований простір $E$ --- комплексний. Спершу кілька деталей.\\
Нехайй $E$ --- комплексний лінійний нормований простір. Розглянемо одночасно $E_{\mathbb{R}}$ --- асоційований з $E$ дійсний нормований простір; тобто під час множення на скаляр ми допускаємо лише дійсні коефіцієнти. Зауважимо, що $E_{\mathbb{R}} = E$ як множини, утім не як простори.\\
Розглянемо довільний функціонал $l \colon E \to \mathbb{C}$. Раз $l(x) \in \mathbb{C}$, то для кожного $x \in E$ можна записати функціонал як $l(x) = m(x) + i n(x)$. У цьому випадку $m(x) = \Re l(x),\ n(x) = \Im l(x)$.

\begin{proposition}
Нехай $l \colon E \to \mathbb{C}$ --- лінійний та обмежений функціонал. Тоді $m,n \colon E_{\mathbb{R}} \to \mathbb{R}$ задають лінійний обмежений функціонал.
\end{proposition}

\begin{proof}
Нехай $\alpha,\beta \in \mathbb{R}$ та $x,y \in E$. Тоді ми отримаємо наступне:\\
$l(\alpha x + \beta y) = m(\alpha x + \beta y) + i n(\alpha x + \beta y)$ (з одного боку)\\
$l(\alpha x + \beta y) = \alpha l(x) + \beta l(y) = \alpha (m(x) + i n(x)) + \beta (m(y) + i n(y)) = (\alpha m(x) + \beta m(y)) + i(\alpha n(x) + \beta n(y))$ (з іншого боку).\\
Знаючи, що комплексне число рівне тоді й лише тоді, коли дійсні та уявні частини збігаються, отримаємо\\
$m(\alpha x + \beta y) = \alpha m(x) + \beta m(y)$ \qquad $n(\alpha x + \beta y) = \alpha n(x) + \beta n(y)$.\\
Отже, $m,n$ --- лінійний функціонали.\\
Обмеженість $m$ (аналогічно з $n$) випливає з такго ланцюга нерівностей:\\
$|m(x)| \leq |m(x) + in(x)| = |l(x)| \leq \|l\| \|x\|$.
\end{proof}

\begin{proposition}
$n(x) = -m(ix)$.\\
Іншими словами, ми можемо функціонал $l$ відновити повністю, знаючи функціонал $m$.
\end{proposition}

\begin{proof}
$m(ix) + in(ix) = l(ix) = i l(x) = i(m(x) + in(x)) = -n(x) + im(x)$.\\
$\implies n(x) = -m(ix)$.\\
$l(x) = m(x) -im(ix)$.
\end{proof}

Повернімось назад до теореми Гана-Банаха. Доб'ємо її на випадок, коли $E$ --- комплексний нормований простір.
\begin{proof}
Маємо $E \supset G$ --- два комплексних простори та $E_{\mathbb{R}}, G_{\mathbb{R}}$ --- асоційовані простори. Маємо функціонал $l \colon G \to \mathbb{C}$, який визначається дійсним функціоналом $m \colon G_{\mathbb{R}} \to \mathbb{R}$. Оскільки це дійсний функціонал, ми можемо продовжити до $M \colon E_{\mathbb{R}} \to \mathbb{R}$ зі збереженням норми.\\
Покладемо $L(x) = M(x) - i M(ix)$. Неважко буде довести, що $L$ --- комплексний лінійний функціонал. Залишилося довести, що $\|L\| = \|l\|$. Знову ж таки, достатньо довести $\|L\| \leq \|l\|$. Запишемо $L(x) = |L(x)|e^{i \varphi}$, де $\varphi = \arg L(x)$. Тоді\\
$|L(x)| = e^{-i \varphi} L(x) = L(e^{-i \varphi} x) = M(e^{-i \varphi} x) = |M(e^{-i \varphi}x)| \leq \|M\| \|e^{-i \varphi} x\| = \|m\| \|x\| \leq \|l\| \|x\|$.\\
Отже, $\|L\|$. Ми тут юзали той факт, що $L(y) = M(y)$ при $L(y) \in \mathbb{R}$.
\end{proof}

\begin{remark}
Зауважимо, що якщо $G$ --- лінійна множина (але не підпростір), то теорема Гана-Банаха все одно виконується.\\
У цьому випадку $\bar{G}$ буде підпростором $E$. Функціонал $l$ продовжується неперервним чином на $\tilde{G}$, а далі застосовується доведена теорема.
\end{remark}
\fi

\end{document}
